\chapter{Lecture 13: Iterator Pattern, Friends and Static Members}
\begin{center}
\textit{February 15, 2018}
\end{center}

The encapsulated \code{List} has $\Theta(n^2)$ traversal time. We want to bring it
back down to $\Theta(n)$ without exposing the Node class.

\section{Design Patterns}
\begin{definition}
A \textbf{design pattern} is a commonly occurring, well-tested solution to a
recurring design problem: ``if you have \emph{this} situation, then \emph{this}
technique may help''.
\end{definition}

\section{The Iterator Pattern}
We need to keep track of how much of the list we have already traversed,
without exposing Nodes. Create an \code{Iterator} class, which is given access
to the nodes. It
\begin{itemize}
  \item keeps track of where we are in the traversal, and
  \item acts like a pointer into the list.
\end{itemize}

Inspiration --- walking an array with a pointer:
\begin{lstlisting}[style=cpp]
for (int *p = array; p != array + arraySize; ++p) {
    ... *p ...
}
\end{lstlisting}
So the iterator should support \code{!=}, unary prefix \code{++}, and unary
\code{*}. We also need a way to create an iterator at the start, and a way to
tell that we have reached the end.

\begin{lstlisting}[style=cpp]
class List {
    struct Node;
    Node *theList = nullptr;
  public:
    class Iterator {
        Node *curr;
      public:
        explicit Iterator(Node *curr) : curr{curr} {}
        int &operator*() const { return curr->data; }
        Iterator &operator++() {          // prefix ++
            curr = curr->next;
            return *this;
        }
        bool operator!=(const Iterator &other) const {
            return curr != other.curr;
        }
    };  // end of Iterator

    Iterator begin() { return Iterator{theList}; }
    Iterator end() { return Iterator{nullptr}; }
    void addToFront(int n);
    ...
};  // end of List
\end{lstlisting}
(The methods of \code{Iterator} touch \code{curr->data}, so in practice they
are defined in \code{list.cc}, after \code{List::Node} is defined.)

\begin{lstlisting}[style=cpp]
List l;
l.addToFront(1);
l.addToFront(2);
l.addToFront(3);
for (List::Iterator it = l.begin(); it != l.end(); ++it) {
    cout << *it << endl;     // prints 3, 2, 1 -- in Theta(n) total
}
\end{lstlisting}

\subsection{\code{auto}}
As of C++11, \code{auto x = y;} declares \code{x} with the same type as
\code{y} (automatic type deduction). So the loop header can be
\code{for (auto it = l.begin(); it != l.end(); ++it)}.

\subsection{Range-Based For Loops}
C++11 has built-in support for the Iterator pattern:
\begin{lstlisting}[style=cpp]
for (auto n : l) {        // n is a copy of each element (type of *it)
    cout << n << endl;
}
for (auto &n : l) {       // n is a reference: we can modify the list
    n += 1;
}
\end{lstlisting}
This works for any class \code{C} such that
\begin{enumerate}
  \item \code{C} has methods \code{begin()} and \code{end()} that return objects
        of some class \code{D}, and
  \item \code{D} has \code{operator*}, prefix \code{operator++} and
        \code{operator!=}.
\end{enumerate}
The range-based for loop is just syntactic sugar for the explicit loop above;
it makes the programmer's job easier.

\section{Friends}
We want client code to get iterators only through \code{begin()} and
\code{end()} (otherwise a client could do \code{List::Iterator\{nullptr\}} or
worse). So the Iterator ctor should be private. But then \code{List} itself
could not call it either.

The \code{Iterator} class can grant its private members to \code{List} by
declaring \code{List} a \textbf{friend}:
\begin{lstlisting}[style=cpp]
class List {
    struct Node;
    Node *theList = nullptr;
  public:
    class Iterator {
        Node *curr;
        explicit Iterator(Node *curr) : curr{curr} {}  // now private
      public:
        int &operator*() const;     // by reference: clients can read/write
        Iterator &operator++();     // returns Iterator&: ++ is an update
        bool operator!=(const Iterator &other) const;
        friend class List;          // List may call the private ctor
    };
    Iterator begin() { return Iterator{theList}; }
    Iterator end() { return Iterator{nullptr}; }
    ...
};
\end{lstlisting}

\begin{important}[Friendship]
Friendship breaks encapsulation: a friend can touch all private members. Have
as few friends as possible.
\end{important}

\section{Accessors and Mutators}
\begin{important}[Advice]
Keep fields private.
\end{important}
To provide controlled read/write access, use \textbf{accessors} (getters) and
\textbf{mutators} (setters). A mutator can enforce invariants.
\begin{lstlisting}[style=cpp]
class Vec {
    int x, y;
  public:
    int getX() const { return x; }
    int getY() const { return y; }
    void setX(int v) { x = v; }
    void setY(int v) { y = v; }
};
\end{lstlisting}

\subsection{I/O Operators and Friends}
I/O operators are standalone functions, so they can't see private fields. If
there are no suitable accessors, we can make the function a friend:
\begin{lstlisting}[style=cpp]
class Vec {
    int x, y;
  public:
    ...
    friend std::ostream &operator<<(std::ostream &out, const Vec &v);
};

std::ostream &operator<<(std::ostream &out, const Vec &v) {
    out << v.x << " " << v.y;
    return out;
}
\end{lstlisting}
Note that this does \emph{not} make \code{operator<<} a method; it is still a
standalone function, just one with access to Vec's privates.

\section{Static Fields and Static Methods}
\begin{definition}
A \textbf{static field} is associated with the class itself, not with
individual objects. Think of it as a per-class global variable: there is exactly
one copy, shared by all objects.
\end{definition}
\begin{lstlisting}[style=cpp]
// student.h
struct Student {
    static int objCount;
    int assns, mt, final;
    Student(int assns, int mt, int final)
        : assns{assns}, mt{mt}, final{final} {
        ++objCount;
    }
};
\end{lstlisting}
\begin{important}[C++ rule]
A static field must be \emph{defined} outside the class, in exactly one
\code{.cc} file:
\begin{lstlisting}[style=cpp]
// student.cc
int Student::objCount = 0;
\end{lstlisting}
\end{important}

\subsection{Static Member Functions}
\begin{itemize}
  \item A static member function is associated with the class, not an object.
  \item It does not need an object to be called on, and so it has no
        \code{this} pointer.
  \item Consequently it can only access static fields and call other static
        member functions.
\end{itemize}
\begin{lstlisting}[style=cpp]
struct Student {
    static int objCount;
    static void printObjCount() {
        cout << objCount << endl;
    }
    ...
};

Student billy{60, 70, 80}, bobby{70, 80, 90};
Student::printObjCount();   // outputs 2
\end{lstlisting}

\section*{Summary}
\begin{itemize}
  \item A design pattern is a reusable solution to a common design problem.
  \item Iterator: a nested class that remembers a position and offers
        \code{*}, \code{++}, \code{!=}; \code{begin()}/\code{end()} create
        them; range-based for uses exactly this protocol.
  \item \code{friend} grants access to private members; use sparingly.
  \item Keep fields private; use accessors/mutators.
  \item Static fields/methods belong to the class; define static fields in a
        \code{.cc} file; static methods have no \code{this}.
\end{itemize}
